\documentclass[11pt]{article}
\usepackage[letterpaper,margin=0.83in,headheight=16pt]{geometry}
\usepackage[T1]{fontenc}
\usepackage[utf8]{inputenc}
\usepackage{lmodern}
\usepackage{microtype}
\usepackage{amsmath,amssymb}
\usepackage{xcolor}
\usepackage{hyperref}
\usepackage{enumitem}
\usepackage{booktabs,longtable,tabularx,array}
\usepackage{fancyhdr}
\usepackage{titlesec}
\usepackage{listings}
\usepackage{tikz}
\usepackage{float}
\definecolor{suiteblue}{HTML}{17365D}
\definecolor{suitelight}{HTML}{EAF1F8}
\definecolor{suitegray}{HTML}{556070}
\hypersetup{colorlinks=true,linkcolor=suiteblue,urlcolor=suiteblue,breaklinks=true}
\titleformat{\section}{\Large\bfseries\color{suiteblue}}{\thesection}{0.55em}{}
\titleformat{\subsection}{\large\bfseries\color{suiteblue}}{\thesubsection}{0.55em}{}
\setlist[itemize]{itemsep=0.25em,topsep=0.35em,leftmargin=1.4em}
\setlist[enumerate]{itemsep=0.25em,topsep=0.35em,leftmargin=1.6em}
\setlength{\parindent}{0pt}
\setlength{\parskip}{0.47em}
\emergencystretch=2em
\lstset{basicstyle=\ttfamily\footnotesize,breaklines=true,frame=single,columns=fullflexible,showstringspaces=false}
\pagestyle{fancy}\fancyhf{}
\lhead{Game Design \& Engineering}\rhead{Technical Reference}\cfoot{\thepage}
\begin{document}
{\huge\bfseries\color{suiteblue} Game Design Concepts And Lenses Reference}\par
\vspace{0.15in}
{\normalsize\color{suitegray} Original engineering synthesis; Schell-aligned concepts explicitly identified}\par
\vspace{0.15in}\hrule\vspace{0.15in}

\section{Source and engineering boundary}
Schell (third edition, 2020 copyright) supplies the conceptual design framing. Formal notation, software architectures, validation gates, metrics, consent protocols, and recommended tooling are supplementary engineering practices introduced by this suite. These are not quotations or source-prescribed standards.


\section{Design as a decision system}
A design decision is a hypothesis about what an experience will become when a player acts in a rule-governed system. Capture decision, expected effect, alternatives, confidence, and a test that could change the decision.


\subsection{Operational review}
\begin{enumerate}
\item What decision or player-facing problem is this technique meant to resolve?
\item Which inputs, assumptions, and environmental constraints determine validity?
\item What test or observation could falsify the intended outcome?
\item What negative or unintended consequences should reviewers document?
\end{enumerate}

\section{Elemental tetrad}
Describe mechanics (what changes), story (meaning and narrative), aesthetics (perception and presentation), and technology (enabling substrate). Evaluate changes jointly: a lower input latency may improve control feel but force different animation or camera work.


\subsection{Operational review}
\begin{enumerate}
\item What decision or player-facing problem is this technique meant to resolve?
\item Which inputs, assumptions, and environmental constraints determine validity?
\item What test or observation could falsify the intended outcome?
\item What negative or unintended consequences should reviewers document?
\end{enumerate}

\section{Essential experience}
Write the intended player perception using verbs and felt consequences. Avoid implementation nouns in the first draft. Translate into a short, falsifiable observation statement.


\subsection{Operational review}
\begin{enumerate}
\item What decision or player-facing problem is this technique meant to resolve?
\item Which inputs, assumptions, and environmental constraints determine validity?
\item What test or observation could falsify the intended outcome?
\item What negative or unintended consequences should reviewers document?
\end{enumerate}

\section{Player motivation}
Distinguish intrinsic engagement from contingent rewards. Record what the player chooses without coercion, which constraints support mastery, and what reward loops risk excessive use.


\subsection{Operational review}
\begin{enumerate}
\item What decision or player-facing problem is this technique meant to resolve?
\item Which inputs, assumptions, and environmental constraints determine validity?
\item What test or observation could falsify the intended outcome?
\item What negative or unintended consequences should reviewers document?
\end{enumerate}

\section{Design filters and iteration}
Prioritize tests of fatal risk before polish. Each prototype answers a single question: what evidence would tell us to keep, change, or cancel the concept?


\subsection{Operational review}
\begin{enumerate}
\item What decision or player-facing problem is this technique meant to resolve?
\item Which inputs, assumptions, and environmental constraints determine validity?
\item What test or observation could falsify the intended outcome?
\item What negative or unintended consequences should reviewers document?
\end{enumerate}

\section{Lens-based design review}
The lenses form a set of alternative questions, not an automated scoring model. Choose lenses for the design uncertainty under review; record why each was relevant and any conflicting conclusions.


\subsection{Operational review}
\begin{enumerate}
\item What decision or player-facing problem is this technique meant to resolve?
\item Which inputs, assumptions, and environmental constraints determine validity?
\item What test or observation could falsify the intended outcome?
\item What negative or unintended consequences should reviewers document?
\end{enumerate}

\section{Interest curves and player attention}
An interest curve is a design/observation tool, not a universal physiological signal. Use annotated timestamps, self-reports, and observable behavior; triangulate rather than claim direct mind-reading.


\subsection{Operational review}
\begin{enumerate}
\item What decision or player-facing problem is this technique meant to resolve?
\item Which inputs, assumptions, and environmental constraints determine validity?
\item What test or observation could falsify the intended outcome?
\item What negative or unintended consequences should reviewers document?
\end{enumerate}

\section{Design ethics}
Consider impacts on players, bystanders, communities, vulnerable populations, and client stakeholders. Revenue and session duration alone do not establish player benefit.


\subsection{Operational review}
\begin{enumerate}
\item What decision or player-facing problem is this technique meant to resolve?
\item Which inputs, assumptions, and environmental constraints determine validity?
\item What test or observation could falsify the intended outcome?
\item What negative or unintended consequences should reviewers document?
\end{enumerate}

\section{Reference evidence register}
\begin{tabularx}{\linewidth}{@{}p{0.13\linewidth}X X@{}}\toprule ID & Evidence or design decision & Verification method\\\midrule
R-01 & \rule{0.95\linewidth}{0.4pt} & \rule{0.95\linewidth}{0.4pt}\\[0.4cm]
R-02 & \rule{0.95\linewidth}{0.4pt} & \rule{0.95\linewidth}{0.4pt}\\[0.4cm]
R-03 & \rule{0.95\linewidth}{0.4pt} & \rule{0.95\linewidth}{0.4pt}\\\bottomrule\end{tabularx}

\end{document}
